% %%
% %% The abstract is a short summary of the work to be presented in the
% %% article.
% \begin{abstract}
% Text-guided image editing aims to modify real images according to natural language instructions while preserving non-target visual content. However, existing diffusion- and rectified flow-based methods often face a trade-off between editing effectiveness and source fidelity, due to inversion errors and unintended global changes. We present a \textbf{training-free text-guided image editing framework based on Visual Autoregressive Modeling (VAR)}, which edits real images by operating directly on hierarchical multi-scale token generation, without reconstructing latent trajectories. Given a source image, our method performs dual-path source encoding: continuous VAE features guide source generation, while discrete image tokens are used for token replacement. We explicitly construct attention-based masks across scales, including a source focus mask, a low-attention preserve mask, and a source-target union mask, to determine which regions should be updated and which should remain anchored to source tokens. By selectively rewriting only mask-targeted tokens and preserving source tokens in non-target regions, the framework enables fine-grained localized edits while maintaining structural consistency and background fidelity. Experiments on standard benchmarks show that our method achieves editing quality and source preservation comparable to recent training-free baselines, demonstrating stable and competitive performance.
% \end{abstract}

% %%
% %% The code below is generated by the tool at http://dl.acm.org/ccs.cfm.
% %% Please copy and paste the code instead of the example below.
% %%
% \begin{CCSXML}
% <ccs2012>
%     <concept>
%         <concept_id>10010147.10010178.10010224.10010225</concept_id>
%         <concept_desc>Computing methodologies~Computer vision tasks</concept_desc>
%         <concept_significance>500</concept_significance>
%     </concept>
%     <concept>
%         <concept_id>10010147.10010371.10010382.10010383</concept_id>
%         <concept_desc>Computing methodologies~Image processing</concept_desc>
%         <concept_significance>500</concept_significance>
%     </concept>
%     <concept>
%         <concept_id>10010147.10010257.10010293</concept_id>
%         <concept_desc>Computing methodologies~Machine learning approaches</concept_desc>
%         <concept_significance>100</concept_significance>
%     </concept>
%     <concept>
%         <concept_id>10010147.10010178.10010224.10010245.10010254</concept_id>
%         <concept_desc>Computing methodologies~Reconstruction</concept_desc>
%         <concept_significance>300</concept_significance>
%     </concept>
% </ccs2012>
% \end{CCSXML}

% \ccsdesc[500]{Computing methodologies~Computer vision tasks}
% \ccsdesc[500]{Computing methodologies~Image processing}
% \ccsdesc[100]{Computing methodologies~Machine learning approaches}
% \ccsdesc[300]{Computing methodologies~Reconstruction}
% %%
% %% Keywords. The author(s) should pick words that accurately describe
% %% the work being presented. Separate the keywords with commas.
% \keywords{training-free image editing, text-guided image editing, visual autoregressive modeling, image inversion, structure preservation}


%%
%% The abstract is a short summary of the work to be presented in the
%% article.
\begin{abstract}
Text-guided image editing aims to modify real images according to natural language instructions while preserving non-target visual content. Existing diffusion- and rectified flow-based methods face a fundamental trade-off between editing effectiveness and source fidelity, stemming from inversion errors inherent to continuous latent trajectories and cross-modal feature entanglement that causes unintended global changes. We present a \textbf{training-free text-guided image editing framework based on Visual Autoregressive Modeling (VAR)}, which sidesteps both challenges by operating on discrete, scale-wise token generation rather than continuous noise removal. Since no trajectory exists to invert, tokens can be selectively replaced at the sampling stage without altering the global generation process. Our method leverages cross-modal attention maps between text and visual tokens to automatically derive high-attention focus masks and low-attention preserve masks at each generation scale — eliminating the need for external segmentation tools or per-category hyperparameter tuning. For real image editing, the source image is encoded into discrete per-scale bitwise tokens via BSQ-VAE, enabling direct token-level background substitution that faithfully preserves source content without inversion or continuous feature injection. Guided by an \textit{understanding-before-editing} philosophy, the model first resolves which spatial regions correspond to the edit via its own cross-modal attention, then rewrites only the targeted tokens. Experiments on standard benchmarks demonstrate that our method achieves superior source fidelity and competitive editing quality over recent training-free baselines, while providing a simpler, faster, and inversion-free pipeline.
\end{abstract}

%%
%% The code below is generated by the tool at http://dl.acm.org/ccs.cfm.
%% Please copy and paste the code instead of the example below.
%%
\begin{CCSXML}
<ccs2012>
    <concept>
        <concept_id>10010147.10010178.10010224.10010225</concept_id>
        <concept_desc>Computing methodologies~Computer vision tasks</concept_desc>
        <concept_significance>500</concept_significance>
    </concept>
    <concept>
        <concept_id>10010147.10010371.10010382.10010383</concept_id>
        <concept_desc>Computing methodologies~Image processing</concept_desc>
        <concept_significance>500</concept_significance>
    </concept>
    <concept>
        <concept_id>10010147.10010257.10010293</concept_id>
        <concept_desc>Computing methodologies~Machine learning approaches</concept_desc>
        <concept_significance>100</concept_significance>
    </concept>
    <concept>
        <concept_id>10010147.10010178.10010224.10010245.10010254</concept_id>
        <concept_desc>Computing methodologies~Reconstruction</concept_desc>
        <concept_significance>300</concept_significance>
    </concept>
</ccs2012>
\end{CCSXML}

\ccsdesc[500]{Computing methodologies~Computer vision tasks}
\ccsdesc[500]{Computing methodologies~Image processing}
\ccsdesc[100]{Computing methodologies~Machine learning approaches}
\ccsdesc[300]{Computing methodologies~Reconstruction}
%%
%% Keywords. The author(s) should pick words that accurately describe
%% the work being presented. Separate the keywords with commas.
\keywords{training-free image editing, text-guided image editing, visual autoregressive modeling, inversion-free editing, cross-modal attention}


